A bank that rates its obligors on a master scale must say, for each grade, how often a rating of that grade is followed by default within a year. When the grades come from an agency or when an internal system is validated against one, that probability of default (PD) is read from a migration matrix: the ratings held at the start of a calendar year are counted by their state at its end, and the share that ended in default is the PD of the grade \citep{anderson1957,jarrow1997,israel2001,jafry2004}. The long-run average of those annual rates is the quantity the European Banking Authority asks a bank to estimate for a through-the-cycle grade \citep{eba2017}, and the one-year matrix is chained to obtain lifetime PDs and stress-test projections \citep{engelmann2024}. A validator who must decide whether a grade's PD is correct therefore depends on how the count handles a rating that leaves the cohort during the year without defaulting in view: the rating is \wdn{} by the agency, because the issuer repaid, stopped paying for the rating, or went private.

A \wdn{} rating hides its outcome. Either it would have defaulted within the year, in which case the PD of its grade is larger than the observed defaults suggest, or it would not. Three treatments are in use. The transition page of the central repository of the European Securities and Markets Authority removes the \wdn{} ratings from the cohort and renormalises, the treatment that \citet{schuermann2004} attribute to industry practice; keeping them in the denominator without following them counts each as a survivor; and the guideline of the European Banking Authority asks to follow them to the year end \citep{eba2017}. Each is right under one assumption about why ratings are \wdn{}. Removal is right when withdrawal is independent of default, keeping is right when no \wdn{} rating would have defaulted, and following is right when the year-end state of every \wdn{} rating can be observed, which a published migration matrix does not allow. When the odds of withdrawal are higher for a rating that is about to default, all three understate the PD, and no test of the matrix sees it, because the observed counts of the cohort that remains are consistent with a Markov chain that is stable from year to year.

The size of the problem is not small. In the generator used here, built from the migration counts of S\&P on a seven-grade scale, a withdrawal ten times as likely in the year before default makes the removed estimate of grade B's PD \nv{abias.fy.i9.industry}\% too low in the first annual cohort and \nv{abias.pool.i9.industry}\% too low when five annual cohorts are pooled; keeping the \wdn{} ratings in the denominator gives \nv{abias.fy.i9.unfollowed}\% and \nv{abias.pool.i9.unfollowed}\%. On the agency's own data the two pages of the repository that count defaults disagree by more than that. For S\&P's grade B, the transition page counts \nv{real.sp.B.D} defaulted ratings in the pooled cohorts and the default page counts \nv{real.sp.B.A}; over the years in which both pages hold data the transition page's annual rate is \nv{real.sp.B.d4d2}\% of the default page's. Part of that gap is definition and part may be informative withdrawal; the data cannot say which.

The proposed method does not choose a treatment. It states the PD of a grade as a function $\mathrm{PD}(k)$ of the odds ratio $k$ of withdrawal for ratings that default within the year, so that $k=1$ is the usual estimate, $k=0$ keeps the \wdn{} ratings as survivors and $k\to\infty$ is the worst case of \citet{manski1989}. The family is the delta-adjusted tipping-point analysis that clinical trials use for outcomes missing not at random \citep{cro2020,carpenter2013}, with the odds ratio as the sensitivity parameter. A decision taken at a threshold, such as a master-scale boundary or a step of the risk-weight function, has a tipping $k^{*}$ at which it flips; a decision that flips only at an implausible $k^{*}$ can be taken as it stands. Where a second source counts the defaults of the same grade differently, the method reads $k$ from it: this is the anchor, and the sampling interval of the second source, carried through the family, gives an anchored interval for the PD.

The contributions are the following.
\begin{itemize}
\item The PD under informative withdrawal as a closed-form family with a closed-form inverse, a closed-form tipping $k$ and a closed-form anchored $k$, with the properties that make $k$ a reparameterisation of the identified set rather than an extra assumption (Sections~\ref{sec:proposed-method} and~\ref{sec:theory}).
\item A proof that, when the second source counts exactly the defaults the \wdn{} ratings would have had, the anchored interval contains the Jeffreys interval of the oracle that follows every withdrawal, so its coverage cannot be below that oracle's; and a derivation of the bias that a second source with a different definition imposes, with the capture at which anchoring stops beating removal.
\item A known-truth evaluation with \nv{reps} repetitions per cell on rating paths generated from the migration counts of S\&P, comparing the removed, kept and followed treatments, a continuous-time estimator that reads withdrawal as independent censoring, the Manski band, the anchored band alone and the anchored interval, over informativeness, cohort size, and second sources that over- or under-count (Sections~\ref{sec:experimental-design} to~\ref{sec:ablation}).
\item An application to the agencies' real cohorts of the central repository, with the anchored $k$ by year and by grade, its stability across years, and the tipping $k$ of the step to the next grade, together with a measurement on known truth of how stable the anchored $k$ is when the two sources define default alike and when their definitions differ.
\end{itemize}

The main results are of two kinds. On known truth the anchored interval covers grade B's true PD of \nv{truthB}\% in at least \nv{h1.min}\% of the \nv{reps} repetitions at informativeness 1, 3 and 9 in the first-year and pooled designs and at 9 in the ten-fold design, against \nv{cov.fy.i9.industry}\% for the removed estimate in the first annual cohort and \nv{cov.pool.i9.industry}\% in the pooled cohorts; the price is a width of \nv{w.fy.i9.anchored} percentage points in the first annual cohort, against \nv{w.fy.i9.followed} for the interval of an oracle that follows every withdrawal. On the agencies' cohorts the pooled anchored $k$ of S\&P's CCC to C grade is \nv{real.sp.CCC-C.k}, with a band from \nv{real.sp.CCC-C.klo} to \nv{real.sp.CCC-C.khi}, but the annual cohorts' own bands hold the pooled value in only \nv{real.sp.B.share}\% of the years for grade B, whereas on known truth with a matching second source they hold it in at least \nv{stab.i9.c1.B.share}\%. The anchored $k$ on agency data is therefore conditional on the two sources defining default alike, which the repository does not guarantee. The paper reports both outcomes as they are.

The rest of the paper is organised as follows. Section~\ref{sec:related-work} places the method among the treatments of \wdn{} ratings, the censoring and competing-risk literature, the sensitivity analyses of missing data and partial identification. Section~\ref{sec:problem} states the problem and its assumptions, Section~\ref{sec:proposed-method} the method and Section~\ref{sec:theory} its properties. Sections~\ref{sec:experimental-design} to~\ref{sec:ablation} give the design of the experiments, their results and the ablation, and Sections~\ref{sec:discussion} to~\ref{sec:conclusions} discuss them, state the limitations, the reproduction commands and the conclusions.
